Differentiable partial differential equation (PDE) solvers underpin
solver-in-the-loop ML training, gradient-based optimal control, and
inverse problems, yet the practical cost of obtaining correct, usable
gradients from a given solver on a given problem is largely
undocumented. Integration effort, computational cost, gradient accuracy,
and numerical conditioning vary widely across solvers and are
discoverable only by trial and error.
%
We introduce \textit{Mosaic}, an extensible benchmarking framework for
differentiable PDE solvers that standardizes access to solver gradients.
Each solver is packaged as a containerized component
(\textit{Tesseract}) exposing a uniform gradient API regardless of
language or automatic differentiation (AD) strategy, enabling
researchers to evaluate, compare, and build on non-trivial physical solvers.
%
Our evaluation of 14 solvers across fluid dynamics, structural
mechanics, and heat transfer demonstrates that the benchmark surfaces
practically relevant differences: order-of-magnitude variation in
computational cost and Jacobian conditioning, alongside setup
incompatibilities and instabilities that rule out solvers for specific
tasks. Despite this variation, all solvers that
produce gradients converge to similar optima, indicating that the
practical barriers are memory limits, numerical stability, and setup
compatibility rather than gradient accuracy alone.
%
Mosaic is open-source and available at \url{https://github.com/pasteurlabs/mosaic}.
